\section{Rule Sets}

While developing rule-based transformations for conjure oxide, it is useful to understand the structure of the rulesets and the `types' of rules that can be used in conjure oxide. Let us first look at how rules actually function, not programmatically, but in an abstract sense. An understanding of (and some experience of) functional programming is incredibly helpful\footnote{For this, I am partial to Graham Hutton's excellent \emph{Programming in Haskell (2nd Ed)}. Also take a look at \emph{Learn You a Haskell for Great Good!} by Miran Lipovača, especially if you prefer a lighter, more witty read.}. Also useful is an understanding of the idea behind graph machines\footnote{Take a look at this blogpost https://amelia.how/posts/the-gmachine-in-detail.html for a good explanation of graph machines (and bonus on lazy evaluation!)} and an understanding of the difference between function results and side effects\footnote{I first gained an understanding of this from the very source -- Alonzo Church's Paper \emph{The Calculi of Lambda Conversion}, which has a very thorough explanation of the concept. However, for a far more accessible expanation is that a Function is a mapping from some input to some output. Anything that the function `permanently' changes in the world outside its scope is a side effect}

In conjure oxide, the rules are functions that take in the expression tree and the symbol table as arguments and return a function result \footnote{\texttt{std::Result<T, E>} in rust records either a Success (\texttt{Ok(T)}) or a Failure (\texttt{Err(E)}). This type of 'error reading' is a functional concept -- much like the \texttt{Either} Monad in Haskell. This is different from throwing exceptions like in imperative languages}, meaning that the original expression tree and symbol table are only modified through side effects of `rule' functions\footnote{Worth noting that this is particularly nice in rust because of the ownership system and the error system}.

There are quite a few advantages to this system, including a few small quality of life things such as the ability to write more descriptive errors and and to pattern match on function results. The most significant, however, is the ability to express rules as functions that can do whatever they need to do as long as they return a failure or a success. This means that an application failure can be treated as a recoverable result rather than a crash. Each function is self contained\footnote{Owing to the Ownership system in Rust, any extra memory allocated in the function is freed as soon as the function exits} meaning that the only things preserved are the initial Symbol Table and the Expression, along with the result being passed along the call stack. The most significant, however, is that this allows for code to be written in a way that enforces that the errors be handled. This is quite useful in general because it means that the codebase can leverage Rust's type system to eliminate edge cases everywhere. The rule aoplication is able to use it to avoid a situation where the rule engine fails unexpectedly and loses a large amount of work, or worse, a situation where a rule seems to have been applied and is visible in the trace but has not actually affected the tree because it failed\footnote{As long as a program doesn't just call unwrap on a rule result... But if you do make this error, you can take some solace in knowing that cloudflare beat you to it!}. 

Now that we know how they are called, we can move on to how the rules themselves are designed. Broadly, there are two categories that we can divide the rules into: Transformation Rules and Representation rules. This is because despite being applied in the same `step' in the process, they have different purposes from each other, but all of the rules of either type all share a common purpose. 

\subsection{Representation Rules}

Essence (which is the input language of conjure-oxide) defines domains which are not present in the type systems of the different solvers' input languages, meaning they need to be encoded in some way into the input languages. The encoding is done using \textit{representation rules}. They are a `type' of rule which implements shared behaviour using a trait\footnote{For details, check out the documentation}. The representation rules must change essence expressions into the target solver's input languiage while preserving all of the relevant information. 